\documentclass{article}

\def\npart {IA}
\def\nterm {Lent}
\def\nyear {2024}
\def\nlecturer {Prof. S.\ Hartnoll}
\def\ncourse {Dynamics and Relativity}

\setcounter{tocdepth}{2}

\usepackage{alltt}
\usepackage{amsmath, amsfonts, amssymb, amsthm}
\usepackage{booktabs}
\usepackage{caption}
\usepackage{enumitem}
\usepackage{fancyhdr}
\usepackage{graphicx}
\usepackage{mathdots}
\usepackage{mathtools}
\usepackage{microtype}
\usepackage{multirow}
\usepackage{pdflscape}
\usepackage{pgfplots}
\usepackage{siunitx}
\usepackage{textcomp}
\usepackage{slashed}
\usepackage{tabularx}
\usepackage{tikz}
\usepackage{tkz-euclide}
\usepackage[normalem]{ulem}
\usepackage[all]{xy}

\pgfplotsset{compat=1.18}
\pagestyle{fancyplain}

\lhead{\emph{\nouppercase{\leftmark}}}
\rhead{
  \ifnum\thepage=1
  \else
    \npart\ \vline\ \ncourse
  \fi}

\def\nauthor{Marcus Ng}
\author{Based on lectures by \nlecturer \\\small Notes taken by \nauthor}
\date{\nterm\ \nyear}
\title{Part \npart\ --- \ncourse}

\newcommand*{\Cdot}{{\raisebox{-0.25ex}{\scalebox{1.5}{$\cdot$}}}}
\setlist[enumerate,1]{label={(\roman*)}}

\newcommand {\pd}[2][ ]{
  \ifx #1 { }
    \frac{\partial}{\partial #2}
  \else
    \frac{\partial^{#1}}{\partial #2^{#1}}
  \fi
}

% Theorems
\theoremstyle{definition}
\newtheorem*{aim}{Aim}
\newtheorem*{axiom}{Axiom}
\newtheorem*{claim}{Claim}
\newtheorem*{cor}{Corollary}
\newtheorem*{conjecture}{Conjecture}
\newtheorem*{defi}{Definition}
\newtheorem*{eg}{Example}
\newtheorem*{ex}{Exercise}
\newtheorem*{fact}{Fact}
\newtheorem*{law}{Law}
\newtheorem*{lemma}{Lemma}
\newtheorem*{notation}{Notation}
\newtheorem*{prop}{Proposition}
\newtheorem*{question}{Question}
\newtheorem*{problem}{Problem}
\newtheorem*{rrule}{Rule}
\newtheorem*{thm}{Theorem}
\newtheorem*{assumption}{Assumption}
\newtheorem*{assert}{Assertion}

\newtheorem*{remark}{Remark}
\newtheorem*{warning}{Warning}
\newtheorem*{exercise}{Exercise}

\newtheorem{nthm}{Theorem}[section]
\newtheorem{nlemma}[nthm]{Lemma}
\newtheorem{nprop}[nthm]{Proposition}
\newtheorem{ncor}[nthm]{Corollary}

\SetEnumitemKey{cases}{
  label=\underline{Case~\arabic*:},
  ref={\arabic*},
  align=left,
  leftmargin=\parindent
}

% \renewcommand{\labelitemi}{--}
% \renewcommand{\labelitemii}{$\circ$}
% \renewcommand{\labelenumi}{(\roman{*})}

\let\stdsection\section
\renewcommand\section{\newpage\stdsection}

% Strike through
\def\st{\bgroup \ULdepth=-.55ex \ULset}


%%%%%%%%%%%%%%%%%%%%%%%%%
%%%%% Maths Symbols %%%%%
%%%%%%%%%%%%%%%%%%%%%%%%%

% Logical symbols
\newcommand{\contradiction}{
  \begin{tikzpicture}[rotate=45,x=0.5ex,y=0.5ex]
  \draw[line width=.1ex] (0,2) -- (3,2) (0,1) -- (3,1) (1,3) -- (1,0) (2,3) -- (2,0);
  \end{tikzpicture}
}


\let\U\relax
\let\C\relax
\let\G\relax

% Matrix groups
% \newcommand{\GL}{\mathrm{GL}}
% \newcommand{\Or}{\mathrm{O}}
% \newcommand{\PGL}{\mathrm{PGL}}
% \newcommand{\PSL}{\mathrm{PSL}}
% \newcommand{\PSO}{\mathrm{PSO}}
% \newcommand{\PSU}{\mathrm{PSU}}
% \newcommand{\SL}{\mathrm{SL}}
% \newcommand{\SO}{\mathrm{SO}}
% \newcommand{\Spin}{\mathrm{Spin}}
% \newcommand{\Sp}{\mathrm{Sp}}
% \newcommand{\SU}{\mathrm{SU}}
\newcommand{\U}{\mathrm{U}}
% \newcommand{\Mat}{\mathrm{Mat}}

% Matrix algebras
% \newcommand{\gl}{\mathfrak{gl}}
% \newcommand{\ort}{\mathfrak{o}}
% \newcommand{\so}{\mathfrak{so}}
% \newcommand{\su}{\mathfrak{su}}
% \newcommand{\uu}{\mathfrak{u}}
% \renewcommand{\sl}{\mathfrak{sl}}

% Special sets
\newcommand{\C}{\mathbb{C}}
\newcommand{\CP}{\mathbb{CP}}
\newcommand{\GG}{\mathbb{G}}
\newcommand{\N}{\mathbb{N}}
\newcommand{\Q}{\mathbb{Q}}
\newcommand{\R}{\mathbb{R}}
\newcommand{\RP}{\mathbb{RP}}
\newcommand{\T}{\mathbb{T}}
\newcommand{\Z}{\mathbb{Z}}
\renewcommand{\H}{\mathbb{H}}

% Brackets
\newcommand{\abs}[1]{\left\lvert #1\right\rvert}
\newcommand{\bket}[1]{\left\lvert #1\right\rangle}
\newcommand{\brak}[1]{\left\langle #1 \right\rvert}
\newcommand{\braket}[2]{\left\langle #1\middle\vert #2 \right\rangle}
\newcommand{\bra}{\langle}
\newcommand{\ket}{\rangle}
\newcommand{\norm}[1]{\left\lVert #1\right\rVert}
\newcommand{\normalorder}[1]{\mathop{:}\nolimits\!#1\!\mathop{:}\nolimits}
\newcommand{\tv}[1]{|#1|}
\renewcommand{\vec}[1]{\boldsymbol{\mathbf{#1}}}

% not-math
% \newcommand{\bolds}[1]{{\bfseries #1}}
% \newcommand{\cat}[1]{\mathsf{#1}}
% \newcommand{\ph}{\,\cdot\,}
% \newcommand{\term}[1]{\emph{#1}\index{#1}}
% \newcommand{\phantomeq}{\hphantom{{}={}}}

% Probability
% \DeclareMathOperator{\Bernoulli}{Bernoulli}
% \DeclareMathOperator{\betaD}{beta}
% \DeclareMathOperator{\bias}{bias}
% \DeclareMathOperator{\binomial}{binomial}
% \DeclareMathOperator{\corr}{corr}
% \DeclareMathOperator{\cov}{cov}
% \DeclareMathOperator{\gammaD}{gamma}
% \DeclareMathOperator{\mse}{mse}
% \DeclareMathOperator{\multinomial}{multinomial}
% \DeclareMathOperator{\Poisson}{Poisson}
% \DeclareMathOperator{\var}{var}
% \newcommand{\E}{\mathbb{E}}
% \newcommand{\Prob}{\mathbb{P}}

% Algebra
% \DeclareMathOperator{\adj}{adj}
% \DeclareMathOperator{\Ann}{Ann}
% \DeclareMathOperator{\Aut}{Aut}
% \DeclareMathOperator{\Char}{char}
% \DeclareMathOperator{\disc}{disc}
% \DeclareMathOperator{\dom}{dom}
% \DeclareMathOperator{\fix}{fix}
% \DeclareMathOperator{\Hom}{Hom}
% \DeclareMathOperator{\id}{id}
% \DeclareMathOperator{\image}{image}
% \DeclareMathOperator{\im}{im}
% \DeclareMathOperator{\re}{re}
% \DeclareMathOperator{\tr}{tr}
% \DeclareMathOperator{\Tr}{Tr}
% \newcommand{\Bilin}{\mathrm{Bilin}}
% \newcommand{\Frob}{\mathrm{Frob}}

% Others
% \newcommand\ad{\mathrm{ad}}
% \newcommand\Art{\mathrm{Art}}
% \newcommand{\B}{\mathcal{B}}
% \newcommand{\cU}{\mathcal{U}}
% \newcommand{\Der}{\mathrm{Der}}
% \newcommand{\D}{\mathrm{D}}
% \newcommand{\dR}{\mathrm{dR}}
% \newcommand{\exterior}{\mathchoice{{\textstyle\bigwedge}}{{\bigwedge}}{{\textstyle\wedge}}{{\scriptstyle\wedge}}}
% \newcommand{\F}{\mathbb{F}}
% \newcommand{\G}{\mathcal{G}}
% \newcommand{\Gr}{\mathrm{Gr}}
% \newcommand{\haut}{\mathrm{ht}}
% \newcommand{\Hol}{\mathrm{Hol}}
% \newcommand{\hol}{\mathfrak{hol}}
% \newcommand{\Id}{\mathrm{Id}}
% \newcommand{\lie}[1]{\mathfrak{#1}}
% \newcommand{\op}{\mathrm{op}}
% \newcommand{\Oc}{\mathcal{O}}
% \newcommand{\pr}{\mathrm{pr}}
% \newcommand{\Ps}{\mathcal{P}}
% \newcommand{\pt}{\mathrm{pt}}
% \newcommand{\qeq}{\mathrel{``{=}"}}
% \newcommand{\Rs}{\mathcal{R}}
% \newcommand{\Vect}{\mathrm{Vect}}
% \newcommand{\wsto}{\stackrel{\mathrm{w}^*}{\to}}
% \newcommand{\wt}{\mathrm{wt}}
% \newcommand{\wto}{\stackrel{\mathrm{w}}{\to}}
% \renewcommand{\d}{\mathrm{d}}
% \renewcommand{\P}{\mathbb{P}}
%\renewcommand{\F}{\mathcal{F}}

\let\Im\relax
\let\Re\relax

% \DeclareMathOperator{\area}{area}
% \DeclareMathOperator{\card}{card}
% \DeclareMathOperator{\ccl}{ccl}
% \DeclareMathOperator{\ch}{ch}
% \DeclareMathOperator{\cl}{cl}
% \DeclareMathOperator{\cls}{\overline{\mathrm{span}}}
% \DeclareMathOperator{\coker}{coker}
% \DeclareMathOperator{\conv}{conv}
% \DeclareMathOperator{\cosec}{cosec}
% \DeclareMathOperator{\cosech}{cosech}
% \DeclareMathOperator{\covol}{covol}
% \DeclareMathOperator{\diag}{diag}
% \DeclareMathOperator{\diam}{diam}
% \DeclareMathOperator{\Diff}{Diff}
% \DeclareMathOperator{\End}{End}
% \DeclareMathOperator{\energy}{energy}
% \DeclareMathOperator{\erfc}{erfc}
% \DeclareMathOperator{\erf}{erf}
% \DeclareMathOperator*{\esssup}{ess\,sup}
% \DeclareMathOperator{\ev}{ev}
% \DeclareMathOperator{\Ext}{Ext}
% \DeclareMathOperator{\fst}{fst}
% \DeclareMathOperator{\Fit}{Fit}
% \DeclareMathOperator{\Frac}{Frac}
% \DeclareMathOperator{\Gal}{Gal}
% \DeclareMathOperator{\gr}{gr}
% \DeclareMathOperator{\hcf}{hcf}
\DeclareMathOperator{\Im}{Im}
% \DeclareMathOperator{\Ind}{Ind}
% \DeclareMathOperator{\Int}{Int}
% \DeclareMathOperator{\Isom}{Isom}
% \DeclareMathOperator{\lcm}{lcm}
% \DeclareMathOperator{\length}{length}
% \DeclareMathOperator{\Lie}{Lie}
% \DeclareMathOperator{\like}{like}
% \DeclareMathOperator{\Lk}{Lk}
% \DeclareMathOperator{\Maps}{Maps}
% \DeclareMathOperator{\orb}{orb}
% \DeclareMathOperator{\ord}{ord}
% \DeclareMathOperator{\otp}{otp}
% \DeclareMathOperator{\poly}{poly}
% \DeclareMathOperator{\rank}{rank}
% \DeclareMathOperator{\rel}{rel}
% \DeclareMathOperator{\Rad}{Rad}
\DeclareMathOperator{\Re}{Re}
% \DeclareMathOperator*{\res}{res}
% \DeclareMathOperator{\Res}{Res}
% \DeclareMathOperator{\Ric}{Ric}
% \DeclareMathOperator{\rk}{rk}
% \DeclareMathOperator{\Rees}{Rees}
% \DeclareMathOperator{\Root}{Root}
% \DeclareMathOperator{\sech}{sech}
% \DeclareMathOperator{\sgn}{sgn}
% \DeclareMathOperator{\snd}{snd}
% \DeclareMathOperator{\Spec}{Spec}
% \DeclareMathOperator{\spn}{span}
% \DeclareMathOperator{\stab}{stab}
% \DeclareMathOperator{\St}{St}
% \DeclareMathOperator{\supp}{supp}
% \DeclareMathOperator{\Syl}{Syl}
% \DeclareMathOperator{\Sym}{Sym}
% \DeclareMathOperator{\vol}{vol}
\numberwithin{equation}{section}

\begin{document}
\maketitle{
    \small 
    \noindent\textbf{Basic Concepts}\\
    Space and time, frames of reference, Galilean transformations. Newton's laws. Dimensional analysis.
    Examples of forces, including gravity, friction and Lorentz.\hspace*{\fill}[4]
    
    \vspace{10pt}
    \noindent\textbf{Newtonian dynamics of a single particle}\\
    Equation of motion in Cartesian and plane polar coordinates. Work, conservative forces and potential energy, motion and the shape of the potential energy function;
    stable equilibria and small oscillations; effect of damping.

    \vspace{5pt}
    \noindent Angular velocity, angular momentum, torque

    \vspace{5pt}
    \noindent Orbits: the $u(\theta)$ equation; escape velocity; Kepler's laws; stability of orbits; motion in a repulsive potential (Rutherford Scattering).

    \vspace{5pt}
    \noindent Rotating frames: centrifugal and Coriolis forces.\\
    *Brief discussion of Foucault pendumum.*\hspace*{\fill}[8]

    \vspace{10pt}
    \noindent\textbf{Newtonian dynamics of systems of particle}\\
    Momentum, angular momentum, energy. Motion relative to the centre of mass; the two body problem. \\
    Variable mass problems; the rocket equation.\hspace*{\fill}[2]

    \vspace{10pt}
    \noindent\textbf{Rigid bodies}\\
    Moments of inertia, angular momentum and energy of a rigid body. 
    Parallel axis theorem. Simple examples of motion involving both rotation and translation (e.g. rolling).\hspace*{\fill}[3]

    \vspace{10pt}
    \noindent\textbf{Special Relativity}\\
    The principle of relativity. Relativity and simultaneity. The invariant interval. Lorents transformations in $(1 + 1)$-dimensional spacetime.
    Time dilation and length contraction. The Minkowski metric for $(1 + 1)$-dimensional spacetime.

    \vspace{5pt}
    \noindent Lorentz transformations in $(3 + 1)$-dimensions. 4-vectors and Lorents invariants. Proper time. 4-velocity and 4-momentum.
    Conservation of 4-momentum in particle decay. Collisions. The Newtonian limit. \hspace*{\fill}[7]
}
\tableofcontents

%%%%%%%%%%%%%%%%%%%%%%
% Lecture 1: 18/01/2024
%%%%%%%%%%%%%%%%%%%%%%
\section{Structure of the Newtonian Universe}

Three dimensional space can be endowed with a \emph{cartesian reference frame} (origin and axes),
such that points in space are labelled:
\[\vec x = (x_1, x_2, x_3)\]
\emph{Time} can be labelled with respect to an arbitrary reference time by a real number $t$.
A \emph{point particle} is an idealised object that is completely determined by it's position at a given time: $\vec x(t)$.
For example we could have $\vec x$ trace the trajectory of a tennis ball travelling through space.
It's \emph{velocity} is the vector
\[\vec v \equiv \diff{\vec x}{t} = \dvec{x}\]
From Vector Calculus, we know that this vector is tangent to the trajectory. Recall that in cartesian coordinates, 
\[\diff{\vec x}{t} = \crbkt{\diff {x_3}{t}, \diff {x_2}{t}, \diff {x_3}{t}}\]
The \emph{acceleration} is:
\[\vec a = \ddvec{x} = \dvec{v} = \diff[2]{x}{t}\]
The above structure is \textbf{not enough} to write down laws of physics.

\begin{quote}
    ``Who's reference frame do we use? Note that they can move with respect to each other and hence produce different observations''
\end{quote}

Consider a free particle that does not experience any forces, i.e; ``nothing happens'' to it.
The position of this particle is given by $\vec x(t)$ but which frame (origin and axis) do we use?
There is a conflict as the particle may be at rest in one frame but moving in a complicated way in another frame.

\subsection{Law of inertia}
There exist \emph{inertial frames}, in which a free particle has a constant velocity.

In an inertial frame we may write
\[\dvec v = \ddvec x = 0\]
The law of inertia is a way to state \emph{Newton's First Law}.

\begin{quote}
    ``As far as we know, this is a true fact about the universe, but not an obvious one!''
\end{quote}

\subsection{Galilean relativity principle}
A frame that is related to an inertial frame by a \emph{Galilean transformation} is also an inertial frame,
and all laws of physics are the same in both frames.

\begin{defi}[Galilean transformation]
    The new frame $\vec x'$ is related to the old frame $\vec x$ by 
    \[\vec x' = R \vec x + \vec k + \vec w t\]
    where
    \begin{itemize}
        \item $R$ is an orthogonal $3x3$ matrix ($R^TR = I$)
        \item $\vec k$ is a constant vector
        \item $\vec w$ is a constant velocity
    \end{itemize}
\end{defi}

For example, a mass dropped from the mast of a moving ship lands at the same place as if the ship were at rest.

It is easy to see that if $\ddvec x = 0 \Leftrightarrow \ddvec x' = 0$.
Hence Galilean relativity constrains the types of forces that are possible.

The Galilean transformations form the Galilean group; often supplemented by time translations
\[t' = t + t_0\]
A result of this is that we treat all inertial frames as having the same time, called \emph{absolute time}
Galilean Relativity implies that the Laws of Physics make reference to no special point, direction, time or velocity.
All of these things can only be relative.
Hence, one cannot be ``at rest'', you have to be at rest with respect to something.

However, \textbf{acceleration is not relative}.
If you are accelerating in an inertial frame, you will be accelerating in all other inertial frames with the same magnitude.
This makes clear that acceleration will be the main object that shows up in the laws of physics.

\section{Forces}
Once there is more than one particle in the universe, there will be interactions between these particles.
In Newtonian physics, these interactions are described by forces.

\subsection{Newtons 2\ts{nd} Law}
In an inertial frame,
\[\dvec p = \vec F\]
where the momentum is defined to be
\[\vec p \equiv m \dvec x\]
for a particle will \emph{intertial mass} $m$. 
The mass is an additional property of point particles, it could change over time, but we will assume that it is constant in time unless stated otherwise.

The force $\vec F$ depends on the interaction, but can only depend on $\vec x$ and $\dvec x$.
Hence, Newtons 2\ts{nd} Law is a second order DE for $\vec x(t)$.
Given $\vec x$ and $\dvec x$ at $t = 0$ for all particles, Newton's equations uniquely determines $\vec x(t)$ for all future times.
\end{document}